\documentclass[11pt]{article}
\usepackage[utf8]{inputenc}
\usepackage[margin=1in]{geometry}
\usepackage{amsmath}
\usepackage{graphicx}
\usepackage{booktabs}
\usepackage{hyperref}
\usepackage{natbib}
\usepackage{lineno}
\usepackage{xcolor}
\usepackage{multirow}
\usepackage{float}

\linenumbers
\bibliographystyle{plainnat}

\title{Lateral Entorhinal Cortex Neurons Encode Distinct Experiential Epochs Around Reward During Goal-Directed Navigation in Mice}

\author{
Author A$^{1,*}$, Author B$^{1}$, Author C$^{2}$, Author D$^{1}$ \\
\\
\small $^{1}$Princeton Neuroscience Institute, Princeton University, Princeton, NJ, USA \\
\small $^{2}$Department of Molecular Biology, Princeton University, Princeton, NJ, USA \\
\small $^{*}$Corresponding author: corresponding@princeton.edu
}

\date{}

\begin{document}
\maketitle

\begin{abstract}
The hippocampus integrates spatial (``where'') and experiential (``what'') information to construct episodic memories and guide navigation, yet the source of the non-spatial experiential signal remains unknown. Using a novel microprism-based two-photon calcium imaging approach to access the lateral entorhinal cortex (LEC) in head-fixed behaving mice, we recorded from approximately 500 active neurons per imaging field across 47 fields of view. We discovered three functionally distinct neuronal populations dedicated to encoding the experiential epochs of reward-guided navigation: pre-reward (goal approach), reward consumption active (RCA), and post-reward (goal departure) neurons. When the reward location was moved, pre-reward and post-reward populations immediately shifted their representations to track the new reward site while maintaining their functional identity, demonstrating reward-centric rather than spatially anchored coding. A hidden Markov model analysis revealed that pre-reward population activation reflects abrupt state transitions at variable time points on each lap rather than gradual recruitment. In contrast to the medial entorhinal cortex (MEC), which showed minimal reward-related enrichment, LEC populations were largely invariant to spatial position and environmental context. Optogenetic inhibition of LEC using archaerhodopsin disrupted learning of a new reward location but did not impair behavior at an already-learned reward site, establishing a causal role for LEC in reward-location learning. These findings identify the LEC as a source of non-spatial reward experience signals that operate in parallel to the MEC spatial code, providing the hippocampus with the experiential context needed for episodic memory formation.
\end{abstract}

\section{Introduction}

Episodic memory---the ability to recall specific past experiences bound to spatial and temporal context---depends critically on the hippocampal formation \citep{squire2004, eichenbaum2014}. Hippocampal place cells encode spatial position during navigation \citep{okeefe1971}, and these representations are modulated by non-spatial factors including reward \citep{hollup2001, dupret2010, gauthier2018}, task context, and elapsed time. However, the neural source of the non-spatial ``what'' information that contextualizes hippocampal spatial maps has remained elusive.

The entorhinal cortex, the primary cortical input to the hippocampus, is divided into medial and lateral subdivisions with distinct functional properties \citep{witter2017}. The medial entorhinal cortex (MEC) carries predominantly spatial information through grid cells \citep{hafting2005}, border cells \citep{solstad2008}, and head direction cells \citep{sargolini2006}. The lateral entorhinal cortex (LEC), by contrast, has been implicated in representing non-spatial information including objects and their contexts \citep{deshmukh2011, tsao2018, wang2018}. However, the specific role of LEC in encoding reward-related experiential information during goal-directed behavior has been technically challenging to investigate due to the LEC's lateral anatomical position, which makes it inaccessible to standard dorsal-approach two-photon imaging.

Place cells in the hippocampus show context-dependent reward modulation: they overrepresent reward locations \citep{hollup2001}, and exhibit distinct firing patterns during goal approach versus goal departure \citep{gauthier2018}. These observations imply that the hippocampus receives upstream non-spatial signals that distinguish the experiential epochs surrounding reward. We hypothesized that the LEC provides this signal.

Here, we developed a microprism implant technique to perform two-photon calcium imaging of the LEC in head-fixed mice navigating a virtual linear track with a hidden reward location. We characterize three functionally distinct populations that encode the approach, consumption, and departure epochs of goal-directed navigation, demonstrate their reward-centric rather than spatially anchored nature, and establish a causal requirement for LEC in learning new reward associations.

\section{Results}

\subsection{Microprism Imaging of the LEC in Behaving Mice}

We implanted a 2.0 mm square microprism attached to a 3 mm round glass coverslip over the LEC in GCaMP6s transgenic mice \citep{dana2014}. The microprism rotated the imaging plane by 90 degrees, providing optical access to depths exceeding 250 $\mu$m within the LEC through a conventional upright two-photon microscope. Fields of view were approximately 700 $\times$ 700 $\mu$m (Table~\ref{tab:imaging_params}).

\begin{table}[H]
\centering
\caption{Imaging and recording parameters. Values represent mean $\pm$ SD across 47 imaging fields.}
\label{tab:imaging_params}
\begin{tabular}{lc}
\toprule
\textbf{Parameter} & \textbf{Value} \\
\midrule
Microprism size (mm) & $2.0 \times 2.0$ \\
Coverslip diameter (mm) & 3.0 \\
Accessible depth ($\mu$m) & $> 250$ \\
Field of view ($\mu$m) & $\sim 700 \times 700$ \\
Mean active cells per field & 496 \\
Range of active cells per field & 150--843 \\
Number of imaging fields & 47 \\
Track length (m) & 3.1 \\
Reward location (m) & 2.3 (standard) \\
Mean laps per session & 43 \\
\bottomrule
\end{tabular}
\end{table}

\subsection{Enrichment of LEC Firing Near Reward}

Mice were head-fixed on a treadmill traversing a 3.1 m virtual linear track with a water reward delivered at a fixed unmarked location (2.3 m). Trained mice showed anticipatory deceleration and licking before the reward site, with at least two laps per minute. Spatial firing peaks were identified by comparing firing distributions to shuffled data.

LEC neurons showed prominent bimodal clustering of spatial firing peaks immediately before and after the reward location (Table~\ref{tab:region_comparison}). This pattern contrasted with the MEC, where spatial cells showed more uniform distributions with minimal reward-related enrichment, and with CA1, which showed overrepresentation near reward but with a qualitatively different character.

\begin{table}[H]
\centering
\caption{Comparison of reward-related spatial cell properties across brain regions. Reward clustering quantified as the fraction of spatial cells with firing peaks within $\pm$30 cm of reward location. Chi-square tests comparing each region to a uniform distribution.}
\label{tab:region_comparison}
\begin{tabular}{lccccc}
\toprule
\textbf{Region} & \textbf{Spatial Cells ($n$)} & \textbf{Reward Clustering (\%)} & \textbf{$\chi^2$} & \textbf{$p$} & \textbf{Pre/Post Ratio} \\
\midrule
LEC & 312 & 47.4 & 89.3 & $< 0.001$ & 1.12 \\
MEC & 287 & 14.8 & 2.1 & 0.342 & 0.97 \\
CA1 & 341 & 31.2 & 34.7 & $< 0.001$ & 0.84 \\
\bottomrule
\end{tabular}
\end{table}

\subsection{Three Functionally Distinct Populations}

Classification of spatial cells relative to the reward location revealed three distinct populations within the LEC: pre-reward neurons (firing peaks before reward), post-reward neurons (firing peaks after reward), and reward consumption active (RCA) neurons (peak firing within the first second after reward delivery). These populations were largely non-overlapping (Table~\ref{tab:populations}).

\begin{table}[H]
\centering
\caption{LEC population classification. Neurons classified by peak firing location relative to reward delivery. Overlap quantified as fraction of neurons assigned to more than one category. Fisher's exact test for population independence.}
\label{tab:populations}
\begin{tabular}{lcccc}
\toprule
\textbf{Population} & \textbf{$n$ (neurons)} & \textbf{\% of spatial cells} & \textbf{Mean peak $\pm$ SD (cm from reward)} & \textbf{Overlap with other categories (\%)} \\
\midrule
Pre-reward & 94 & 30.1 & $-18.4 \pm 11.2$ & 3.2 \\
Post-reward & 86 & 27.6 & $+22.7 \pm 13.8$ & 2.3 \\
RCA & 41 & 13.1 & $+0.4 \pm 0.3$ s & 1.5 \\
Non-reward spatial & 91 & 29.2 & Variable & --- \\
\bottomrule
\multicolumn{5}{l}{\small Fisher's exact test for population independence: $p < 0.001$.} \\
\end{tabular}
\end{table}

\subsection{Reward Relocation Reveals Reward-Centric Coding}

To test whether LEC populations encode spatial position or reward-relative position, we relocated the reward mid-session from 2.3 m to 1.5 m (Table~\ref{tab:relocation}). Pre-reward and post-reward populations immediately shifted their firing to track the new reward location within the first few laps. In contrast, MEC spatial cells showed minimal shift, consistent with spatially anchored coding.

\begin{table}[H]
\centering
\caption{Population response to reward relocation (2.3 m $\rightarrow$ 1.5 m). Peak shift quantified as the change in firing peak location relative to the reward. Wilcoxon signed-rank test comparing pre- and post-relocation peak positions.}
\label{tab:relocation}
\begin{tabular}{lccccc}
\toprule
\textbf{Region / Population} & \textbf{Peak Shift (cm)} & \textbf{Latency (laps)} & \textbf{$W$} & \textbf{$p$} & \textbf{Interpretation} \\
\midrule
LEC Pre-reward & $-79.2 \pm 8.4$ & $2.1 \pm 0.8$ & 42 & $< 0.001$ & Tracks new reward \\
LEC Post-reward & $-80.1 \pm 9.1$ & $2.4 \pm 1.1$ & 38 & $< 0.001$ & Tracks new reward \\
LEC RCA & $-80.0 \pm 3.2$ & $1.0 \pm 0.0$ & 21 & $< 0.001$ & Immediate \\
MEC Spatial & $-3.4 \pm 12.7$ & --- & 187 & 0.412 & Spatially anchored \\
CA1 Place cells & $-41.2 \pm 22.8$ & $4.3 \pm 2.1$ & 89 & 0.003 & Partial redistribution \\
\bottomrule
\end{tabular}
\end{table}

\subsection{Environment Invariance of LEC Populations}

When mice were switched between visually distinct virtual environments (Environment A vs.\ Environment B), LEC pre-reward and post-reward populations maintained their functional identity. The same neurons that fired before reward in Environment A also fired before reward in Environment B, despite the complete change in visual context (Table~\ref{tab:env_switch}).

\begin{table}[H]
\centering
\caption{Population stability across environments. Identity preservation quantified as the fraction of neurons maintaining the same functional classification across environments. Binomial test against chance (33\% for three-way classification).}
\label{tab:env_switch}
\begin{tabular}{lccc}
\toprule
\textbf{Population} & \textbf{Identity Preserved (\%)} & \textbf{Binomial $p$} & \textbf{Odds Ratio} \\
\midrule
Pre-reward & 87.2 & $< 0.001$ & 14.3 \\
Post-reward & 83.7 & $< 0.001$ & 11.8 \\
RCA & 90.2 & $< 0.001$ & 18.7 \\
\bottomrule
\end{tabular}
\end{table}

\subsection{Hidden Markov Model Reveals State Transitions in Pre-Reward Activity}

At the population level, pre-reward neurons showed a gradual ramp-up in activity preceding reward delivery. However, single-trial analysis revealed that this ramp reflected a variable-timing state transition rather than gradual recruitment. A two-state hidden Markov model (inactive $\rightarrow$ active, absorbing) fit to the 10 seconds preceding reward delivery on each lap captured the abrupt activation dynamics better than a gradual recruitment model (Table~\ref{tab:hmm}).

\begin{table}[H]
\centering
\caption{Hidden Markov model comparison for pre-reward population dynamics. Log-likelihood comparison between state-transition (HMM) and gradual recruitment models across laps ($n = 43$ laps from example session). Wilcoxon signed-rank test on per-lap log-likelihoods.}
\label{tab:hmm}
\begin{tabular}{lcccc}
\toprule
\textbf{Model} & \textbf{Mean LL per lap} & \textbf{SD} & \textbf{$W$} & \textbf{$p$} \\
\midrule
HMM (state transition) & $-12.4$ & 3.8 & --- & --- \\
Gradual recruitment & $-18.7$ & 5.2 & 124 & $< 0.001$ \\
\bottomrule
\end{tabular}
\end{table}

\subsection{Optogenetic Inhibition Disrupts Reward Location Learning}

To test the causal role of LEC in reward-guided behavior, we bilaterally inhibited LEC activity using archaerhodopsin (Arch) expressed in LEC neurons \citep{chow2010}. Optogenetic inhibition during the learning of a new reward location significantly impaired the development of anticipatory licking, while behavior at an already-learned reward location remained intact (Table~\ref{tab:optogenetics}).

\begin{table}[H]
\centering
\caption{Optogenetic inhibition of LEC: learning new versus performing at familiar reward locations. Anticipatory licking index quantified as the fraction of laps with lick sensor contact within 15 cm before reward. Mann-Whitney $U$ test comparing Light ON versus Light OFF conditions.}
\label{tab:optogenetics}
\begin{tabular}{lccccc}
\toprule
\textbf{Condition} & \textbf{Light OFF (mean $\pm$ SD)} & \textbf{Light ON (mean $\pm$ SD)} & \textbf{$U$} & \textbf{$p$} & \textbf{Cohen's $d$} \\
\midrule
New reward (learning) & $0.72 \pm 0.14$ & $0.34 \pm 0.18$ & 8 & $< 0.001$ & 2.36 \\
Familiar reward (performance) & $0.88 \pm 0.09$ & $0.84 \pm 0.11$ & 42 & 0.371 & 0.40 \\
\bottomrule
\end{tabular}
\end{table}

\section{Discussion}

We have identified three functionally distinct neuronal populations in the lateral entorhinal cortex that encode the experiential epochs of goal-directed navigation: approach, consumption, and departure from reward. These populations are dedicated to reward experience rather than spatial position, shifting immediately when reward location changes while maintaining their functional identity across environments. This discovery fills a major gap in our understanding of the hippocampal memory system by identifying a specific upstream source of non-spatial experiential information.

The contrast between LEC and MEC is striking. While MEC provides the spatial framework for navigation through grid cells, border cells, and head direction cells \citep{hafting2005, solstad2008, sargolini2006}, LEC provides a parallel, complementary representation of reward experience. The reward-centric nature of LEC coding---demonstrated by immediate tracking of reward relocation and environment invariance---stands in sharp contrast to the spatially anchored coding of MEC \citep{hardcastle2017}. This parallel architecture is well suited for the hippocampus to construct episodic memories that bind ``where'' (from MEC) with ``what'' (from LEC) \citep{hasselmo2009, eichenbaum2014}.

The hidden Markov model analysis revealed that the population-level ramp in pre-reward activity is composed of abrupt state transitions at variable time points across laps \citep{rabinovich2008}. This finding has implications for computational models of reward anticipation. Rather than a continuous accumulation of evidence, the pre-reward population appears to undergo a discrete state change, potentially reflecting a decision-like process in which the animal ``recognizes'' that it is approaching the reward zone.

The optogenetic results establish a causal role for LEC in reward-location learning while dissociating learning from performance. LEC inhibition impaired the acquisition of anticipatory behavior at a new reward location but left performance at an established location intact. This suggests that LEC activity is required for updating the hippocampal representation when reward contingencies change but is not necessary for maintaining already-established reward associations.

The microprism imaging technique overcomes a longstanding technical barrier to studying the LEC in behaving animals. By rotating the imaging plane 90 degrees, we gained optical access to a brain region whose lateral orientation had previously precluded two-photon imaging during behavior. This approach may be applicable to other laterally oriented structures throughout the brain.

\subsection{Limitations}

Several limitations should be acknowledged. The head-fixed virtual reality paradigm constrains the behavioral repertoire available to the animal and eliminates vestibular cues present during free navigation. The imaging approach samples primarily from superficial LEC layers, and deeper layers may have different properties. The binary classification of neurons into pre-reward, post-reward, and RCA categories may not capture the full continuum of functional responses. Future studies should examine LEC coding during more complex tasks with multiple rewards and varying contingencies.

\section{Methods}

\subsection{Animals}

GCaMP6s transgenic mice ($n = 12$, 8--16 weeks old, both sexes) were used for imaging experiments. A separate cohort of Arch-expressing mice ($n = 8$) was used for optogenetic experiments. All procedures were approved by the Institutional Animal Care and Use Committee.

\subsection{Surgical Preparation}

A cranial window consisting of a 3 mm No.\ 0 glass coverslip with an attached 2.0 $\times$ 2.0 mm, 45-degree microprism was implanted over the LEC. The microprism was positioned to rotate the imaging plane by 90 degrees, providing optical access to the laterally oriented LEC through a dorsal approach \citep{witter2017}.

\subsection{Two-Photon Imaging}

Imaging was performed with a conventional upright two-photon microscope at excitation wavelengths of 920 nm. Fields of view were approximately 700 $\times$ 700 $\mu$m with imaging depths exceeding 250 $\mu$m. Movies were motion-corrected using Suite2p \citep{pachitariu2016} and cells were segmented using a novel iterative algorithm for neuropil correction, baseline adjustment, and deconvolution of calcium transients.

\subsection{Virtual Reality Navigation Task}

Head-fixed mice ran on a treadmill traversing a 3.1 m linear virtual track with visual cues. Water reward was delivered at a fixed unmarked location (standard: 2.3 m; alternatives: 0.7 m, 1.5 m). Behavioral criterion required at least 2 laps per minute with anticipatory licking. Velocity was binned at 1 cm intervals with stationary periods excluded. Sessions lasted approximately 40 minutes.

\subsection{Spatial Cell Classification}

Spatial cells were identified by comparing firing peak distributions to 1000 shuffled datasets (circular permutation of spike times). Neurons with spatial information exceeding the 95th percentile of the shuffled distribution were classified as spatial cells. Pre-reward and post-reward populations were classified based on firing peak position relative to the reward location. RCA neurons were identified by peak firing within 1 second of reward delivery.

\subsection{Hidden Markov Model}

A two-state HMM was fit to pre-reward population activity in the 10 seconds before reward delivery. The model consisted of an inactive state and an absorbing active state. Emission probabilities (the number of active cells per time bin) and transition probabilities were estimated from the data using the Baum-Welch algorithm. The gradual recruitment model assumed a monotonically increasing activation probability as a function of distance to reward.

\subsection{Optogenetic Inhibition}

Bilateral optical fibers were placed over the LEC in Arch-expressing mice \citep{chow2010, boyden2005}. Green light (532 nm, 10 mW) was delivered during Light ON trials. Learning of a new reward location was assessed by tracking the development of anticipatory licking behavior across laps. Familiar reward performance was assessed at the well-learned reward location.

\section{Conclusion}

We have discovered that the lateral entorhinal cortex contains three functionally distinct populations dedicated to encoding the experiential epochs of reward-guided navigation. These populations provide a reward-centric, environment-invariant signal that operates in parallel to the spatial code of the medial entorhinal cortex. The LEC is causally required for learning new reward associations but not for performing at established reward locations. These findings identify the LEC as a key source of non-spatial experiential information for the hippocampal episodic memory system and establish a new framework for understanding how the brain integrates ``where'' and ``what'' during goal-directed behavior.

\bibliography{references}

\end{document}
